% Internal Consistency Is Not Evidence — arXiv preprint
% Build: tectonic main.tex  (XeTeX engine; fetches packages automatically)
% STATUS: working draft circulated for co-author review and consent prior to submission.
%
% NOTE ON THE PREAMBLE. This mirrors ../../transit/paper/main.tex, which in
% turn mirrors the standard's paper. `minted` is not loaded: it shells out to
% pygments, which would make the CI paper build depend on a Python package
% nobody installs there. `listings` needs no external tool.
\documentclass[11pt]{article}

\usepackage[margin=1in]{geometry}
\usepackage{tikz}
\usepackage{graphicx}
\usepackage{booktabs}
\usepackage{array}
\usepackage{amsmath, amssymb}
\usepackage{listings}
\usepackage{xcolor}
\usepackage{colortbl}
\usepackage[hidelinks]{hyperref}
\usepackage{enumitem}
\usepackage{microtype}
\usepackage{caption}
\usepackage{siunitx}
\usepackage{amsthm}
\newtheorem{definition}{Definition}
\newtheorem{theorem}{Theorem}
\newtheorem{proposition}{Proposition}
\usepackage[section]{placeins}

% Advanced AI Society brand palette
\definecolor{aailime}{HTML}{CFFF04}
\definecolor{aailimetint}{HTML}{F2FFB8}
\definecolor{aaiink}{HTML}{0A0A0A}
\definecolor{aaired}{HTML}{FFE3E4}
\definecolor{aaineutral}{HTML}{E7E3DD}
\definecolor{aaipaper}{HTML}{F0EDEA}
\definecolor{aaimuted}{HTML}{8F8A82}
\definecolor{aaidarklime}{HTML}{7A9900}
\definecolor{codepanel}{HTML}{F7F5F1}

\lstset{basicstyle=\ttfamily\footnotesize, breaklines=true, frame=single,
        framerule=0.3pt, xleftmargin=1em, columns=fullflexible}
\captionsetup{font=small, labelfont=bf}

\newcommand{\ephemeris}{\texttt{ephemeris}}
\newcommand{\transit}{\texttt{transit}}

% The Society's bracket mark, drawn natively (lime on the dark band)
\newcommand{\aaimark}[1]{%
\begin{tikzpicture}[x=#1, y=#1, line cap=butt]
  \fill[aailime] (0,0) rectangle (0.16,1);
  \fill[aailime] (0.16,0) rectangle (0.42,0.24);
  \fill[aailime] (0.16,0.76) rectangle (0.42,1);
  \fill[aailime] (1.04,0) rectangle (1.20,1);
  \fill[aailime] (0.78,0) rectangle (1.04,0.24);
  \fill[aailime] (0.78,0.76) rectangle (1.04,1);
\end{tikzpicture}}

\begin{document}

% ---------------------------------------------------------------- cover page
\thispagestyle{empty}

\begin{tikzpicture}[remember picture, overlay]
  \fill[aaiink] (current page.north west) rectangle
    ([yshift=-6.35cm]current page.north east);
  \fill[aailime] ([yshift=-6.35cm]current page.north west) rectangle
    ([yshift=-6.53cm]current page.north east);
  \node[anchor=north west, inner sep=0pt]
    at ([xshift=1in, yshift=-0.95cm]current page.north west) {\aaimark{0.62cm}};
  \node[anchor=north west, inner sep=0pt, text width=14.2cm]
    at ([xshift=1in+1.15cm, yshift=-0.98cm]current page.north west)
    {\sffamily\bfseries\footnotesize\color{aailime}%
     ADVANCED AI SOCIETY\; \textcolor{aaimuted}{$\cdot$}\; PROOF-OF-CONTROL
     INITIATIVE\\[3pt]
     {\normalfont\sffamily\footnotesize\color{aaimuted}
      Working draft v0.1 $\cdot$ August 2026 $\cdot$ for co-author review
      prior to submission}};
  \node[anchor=north west, inner sep=0pt, text width=\textwidth]
    at ([xshift=1in, yshift=-2.55cm]current page.north west)
    {\sffamily\bfseries\fontsize{21.5}{26}\selectfont\color{aaipaper}%
     Internal Consistency Is Not Evidence: Custody, Durability, and
     the Limits of a Tamper-Evident Log};
  \node[anchor=north west, inner sep=0pt, text width=\textwidth]
    at ([xshift=1in, yshift=-5.55cm]current page.north west)
    {\sffamily\small\color{aaipaper}%
     \textcolor{aailime}{\rule[0.5ex]{1.6em}{2pt}}\hspace{0.6em}%
     A log that verifies against itself has told you nothing.};
\end{tikzpicture}

\vspace*{5.9cm}

% Compact author block
{\small\noindent
\textbf{Jim Schwoebel}\textsuperscript{1,2}

\smallskip\noindent
{\footnotesize
\textsuperscript{1}\,Quome \quad
\textsuperscript{2}\,Advanced AI Society, \url{https://advancedaisociety.org}\footnote{Prepared
within the Proof-of-Control Initiative of the Advanced AI Society. Authorship is
limited to contributors to this manuscript. Co-authors will be added only upon
explicit consent and substantive contribution; the author list is deliberately
short at this stage rather than inherited from a sibling paper.}}}

\vspace{0.9em}
\noindent\textcolor{aailime}{\rule{\textwidth}{1.6pt}}
\vspace{0.2em}

\begin{center}
{\sffamily\bfseries\small ABSTRACT}
\end{center}
\vspace{-0.5em}
\noindent{\small
A tamper-evident log chains its records, roots them in a Merkle tree, and
signs the root. Break any record and a link breaks. That is all true, and it
is worth almost nothing against the party running the log, because that party
does not break links---they recompute them. The rewritten log is not a damaged
log. It is a well-formed log of different events, and every check it can
perform on itself passes.

We say plainly what does the work instead. It is not cryptography; it is
custody. A root that left the operator's hands before the rewrite is the only
artefact that contradicts the rewrite, and a system with no such publication
point is undefended no matter how carefully its proofs are implemented. We
build the log anyway---chain, RFC~6962 tree, inclusion and consistency proofs
verified exhaustively rather than sampled---and then make the tool print, on
every root it emits, the sentence saying what it does not defend against.

Along the way we hit something we did not expect, and we think it is the more
useful of our two findings. The requirement was fail-closed durability: never
acknowledge a claim whose bytes are not on disk. The obvious way to test that
is to make the store unavailable and check that everything is refused, and the
obvious way to make it unavailable is to delete the file. That test passes
while the system is broken. A process holding an open descriptor keeps writing
after the path is unlinked, \texttt{fsync} returns success, and the log
acknowledges claims backed by bytes nothing can ever read back. The fault
injection injects no fault. We found it by luck---our assertion failed for the
opposite reason---and we report the luck rather than dressing it up as method.

We also report that the schema we were implementing chains over two fields of
a record rather than the record, which leaves the step index, the target
resource, the policy bundle and every identity field outside the chain
entirely.}

\clearpage

% ------------------------------------------------------------------- content

\section{The Problem}
\label{sec:problem}

Somebody hands you a log. It records decisions an automated agent made:
what it wanted to do, what the policy engine said, whether the action went
through. Every record contains a hash of the record before it, so the records
form a chain. The whole set is summarised by a Merkle root, and that root is
signed. You check it. Every link holds. The signature is good.

What have you learned?

The honest answer is: that this log is well-formed. Not that it is true. Those
are different, and the difference is the entire subject of this paper.

Here is why. Suppose the operator of the log wants record three to say
something else. They are not going to alter record three and leave a broken
link for you to find---that would be a strange kind of adversary, one who
tampers but stops halfway. They alter record three, recompute the chain link
for record four, and record five, and so on to the end. Then they sign the new
root. Now the log has no broken link, because they fixed the links. It has no
gap in the sequence, because they did not remove anything. The signature over
the new root verifies, because they hold the key.

You are looking at a log that is, in every mechanical sense, perfect. And it is
a log of events that did not happen.

I want to be careful about what is surprising here, because on one reading it
is not surprising at all. Of course somebody with the signing key can sign
things. That is what a key is. Everybody nods.

The trouble is that the nodding does not survive contact with a deployment. In
practice you get a document that says the log is tamper-evident, and an
auditor who verifies the log and reports that it checks out, and neither the
document nor the auditor makes the distinction I just made. The verification
step is real work, it produces a green tick, and the green tick means
``well-formed,'' and everyone reads it as ``true.'' The gap is not in anybody's
understanding of cryptography. It is in what the tooling reports and what the
report is taken to mean.

So we built the tool, and then we made the tool refuse to produce that green
tick without a sentence next to it.

\subsection{The one thing that does contradict a rewrite}

There is exactly one artefact that catches this, and it is not clever.

If, before the rewrite, the operator published the root---handed it to a
counterparty, wrote it into a transparency log, printed it and put it in a
drawer---then that old root is a commitment they can no longer edit. A
Merkle consistency proof between the old root and the new one either
demonstrates that the new tree extends the old, or it does not exist, and
against a rewritten history it does not exist. RFC~6962 defines exactly this
proof for exactly this reason \cite{rfc6962}.

Notice what that costs. Not computation: the proof is logarithmic and the
verification is a few dozen hashes. What it costs is a relationship. Somebody
other than the operator has to have been holding something, before the fact,
and be willing to say so afterwards. That is a custody arrangement and an
organisational commitment, and no amount of engineering inside the log
produces it.

This is the paper's central claim, and I want to state it in its blunt form
before we dress it up:

\begin{quote}
The security of a tamper-evident log against the party operating it is
bounded above by the custody of its published roots, and is independent of
the quality of its cryptography.
\end{quote}

Independent. Not ``mostly determined by.'' A log with a flawless RFC~6962
implementation and no publication point, and a log with a mediocre one and no
publication point, are equally defenceless against their operator. The
cryptography is doing something---it defends against everyone who is not the
operator, which is most attackers---but against the operator it contributes
zero, and adding more of it continues to contribute zero.

\subsection{Why this matters more for agents than it used to}

Transparency logs have lived with this for a decade and solved it socially:
Certificate Transparency has multiple independent logs, monitors that watch
them, and auditors that gossip roots between clients \cite{rfc6962,rfc9162}.
The design assumes the log operator is not trusted, and builds the ecosystem
that makes the assumption safe.

The systems we are now being asked to log are not like that. An enterprise
deploying an autonomous agent wants a record of what the agent did, for its
own governance and for its regulator. The natural deployment is one log, run
by the enterprise, holding records about the enterprise's own agents. Every
party in that picture is the same party. The evidence that an agent did not
exceed its authority is held, signed, and---if it comes to it---rewritten by
the organisation whose interest is served by the answer.

I do not think this is usually malice. It is more often that nobody drew the
diagram and noticed that the auditor and the audited are the same box.

\section{What Others Have Built}
\label{sec:related}

The machinery here is old and good, and none of it is ours.

Merkle gave us the hash tree \cite{merkle1980}, which lets you commit to a set
of things with one value and later prove membership without revealing the set.
Haber and Stornetta gave us the linked timestamp \cite{haber1991}, and stated
the custody problem more clearly in 1991 than most systems state it now: their
scheme distributes trust across clients precisely because a single timestamping
service could collude with a client to backdate. That is our problem, thirty-five
years earlier, and their answer was structural rather than cryptographic.

Certificate Transparency turned this into deployed infrastructure
\cite{rfc6962,rfc9162}. The two proofs we implement are theirs: inclusion, that
a record is in a tree with a given root; and consistency, that a tree of $n$
leaves extends a tree of $m \le n$ leaves. Crosby and Wallach worked through
what tamper-evident logging costs in practice and where the naive data
structures fall over \cite{crosby2009}.

On the durability side, the write-ahead log and the group commit are from
database work of the seventies and eighties, with ARIES as the canonical
write-up \cite{mohan1992}. The idea we lean on---do not tell anybody the
transaction committed until the bytes are on the disk---is the oldest idea in
this paper.

What is newer, and more directly relevant to Section~\ref{sec:test}, is a line
of work showing that the storage stack does not behave the way applications
believe. Pillai and colleagues found that applications routinely depend on
file-system behaviours that are not guaranteed, and that the resulting bugs
survive testing because the tests share the application's assumptions
\cite{pillai2014}. Rebello and colleagues asked what happens when
\texttt{fsync} itself fails and found that most applications handle it badly
or not at all \cite{rebello2020}. Our finding is a cousin of these: not that
the storage stack lied, but that the \emph{test} of the storage stack was a
no-op for a reason nobody in the room had thought about.

Where we differ from all of this is emphasis rather than mechanism. The
literature above mostly asks how to build the thing correctly. We are asking
what a correctly built thing is worth in a deployment where the operator and
the audited party coincide, and our answer is: less than the documentation
says, and the shortfall is not fixable inside the tool.

\section{Writing It Down Precisely}
\label{sec:formal}

The symbols in this section will not make anything truer than the English
already did. A reader who skips to Section~\ref{sec:test} loses only precision.
But it is worth being exact about which checks depend on what, because the
whole argument is that two of them look alike and are not.

Let a log be a sequence of records $R = (r_0, \dots, r_{n-1})$. Write $d(r)$
for the digest of a record's canonical form---RFC~8785 \cite{rfc8785}, computed
by a shared implementation rather than a second one, for reasons the sibling
paper works through---and define the chain by
$h_0 = 0^{256}$ and $h_{i+1} = H(h_i \| d(r_i))$. Let $\mathrm{MTH}(R)$ be the
RFC~6962 Merkle tree head over the leaf hashes of $R$.

\begin{definition}[Internal consistency]
$R$ is internally consistent if each $r_i$ declares chain head $h_i$, and the
step indices of $R$ are exactly $0, \dots, n-1$.
\end{definition}

\begin{definition}[Extension]
$R'$ extends $R$ if $R$ is a prefix of $R'$.
\end{definition}

Now the two propositions that matter, and the gap between them.

\begin{proposition}[Internal consistency is cheap to forge]
\label{prop:cheap}
For any sequence of claims $C$ and any modification of any claim in $C$, there
is an internally consistent log recording the modified sequence, computable in
time linear in $|C|$ by any party able to sign a root.
\end{proposition}

The proof is the attack: recompute the chain from the modified record forward,
rebuild the tree, sign. Nothing is hard. We ship it as a runnable example
rather than as an argument, because a reader can then watch every internal
check return true.

\begin{proposition}[Consistency against a prior root is not]
\label{prop:prior}
Let $\rho = \mathrm{MTH}(R)$ be published at size $m$. If $R'$ of size
$n \ge m$ does not extend $R$, then no consistency proof between
$(m, \rho)$ and $(n, \mathrm{MTH}(R'))$ verifies, under collision resistance
of $H$.
\end{proposition}

This is RFC~6962's property, not a new one \cite{rfc6962}. We restate it only
to make the asymmetry visible: Proposition~\ref{prop:cheap} needs the signing
key and nothing else, and Proposition~\ref{prop:prior} needs an artefact that
escaped the operator at a particular time. The first is a computation. The
second is a fact about the world.

\begin{quote}
\textbf{The custody condition.} A log detects operator rewriting at size $m$
if and only if some party outside the operator's control holds
$\mathrm{MTH}(R_{\le m})$ and will produce it on demand.
\end{quote}

I have written that as a biconditional and I want to flag that the ``only if''
direction is the interesting one and the one we can actually argue: with no
external holder, the operator's view of history is the only view, and any
internally consistent log is consistent with it. The ``if'' direction assumes
the external party is honest and available, which is an organisational
assumption doing a lot of quiet work. We have not formalised what happens when
the external holder is partly trusted, or when there are several disagreeing
ones. That is where the Certificate Transparency ecosystem put a decade of
effort \cite{rfc9162}, and we have not repeated it.

\section{The Test That Proves Nothing}
\label{sec:test}

Now the part we did not see coming.

Separately from all the Merkle machinery, the log has a durability
requirement, and it is the requirement the whole component exists for: never
acknowledge a claim whose bytes are not durably written. If an enforcement
point is waiting on this acknowledgement before letting an action proceed,
then acknowledging early means an action happens with no evidence that it
happened. Fail-closed, no exceptions, no flag to turn it off.

The implementation is standard. Serialise the batch, one \texttt{write}, one
\texttt{fsync}, and only then construct the acknowledgement. We made the
acknowledgement type constructible in exactly one place in the codebase, after
the sync returns, so that ``acknowledge before durable'' is not an expressible
program rather than a discouraged one.

Then we wrote the test. The requirement says: make the evidence store
unavailable, and observe that actions do not proceed. So: create the log, take
the store away, submit some claims, assert that every one of them is refused.

How do you take the store away? You delete the file. That is the obvious move,
and it is what we did---unlink the segment and put a directory in its place, so
that any attempt to re-open it fails.

The test failed. Every claim was acknowledged.

Here is what is going on, and it is embarrassing in the specific way that good
findings often are. On a POSIX system, a file descriptor refers to an inode,
not to a path. Unlinking the path removes the directory entry. It does not
touch the descriptor. Our committer had opened the segment at start-up and was
still holding that descriptor, so:

\begin{lstlisting}
write to unlinked fd: Ok(())
reopen at same path:  Err(Os { code: 21, kind: IsADirectory })
\end{lstlisting}

The write succeeds. The \texttt{fsync} succeeds. The bytes really are on the
disk, in an inode with no name, which the kernel will free the moment the last
descriptor closes. The durability requirement is satisfied in the narrowest
possible reading---the bytes were synced---and violated in every sense anybody
cares about, because the evidence cannot be read back by anything, ever.

So the log acknowledges the claim. The enforcement point lets the action
through. The action happens, and there is no record of it, and no error was
raised anywhere in the system.

\subsection{Why we want to make a fuss about this}

The bug is not that unlinked descriptors keep working. That is specified
behaviour---\texttt{unlink} removes the directory entry and defers destroying
the file until the last descriptor is closed \cite{posix2018}---and any systems
programmer asked directly will tell you. We confirmed it on macOS with the
standalone probe above and have not run it elsewhere. The finding is about the
test.

We wrote a fault-injection test. It injected no fault. The store was never
unavailable; only its name was. And the assertion the test existed to check
was therefore never exercised at all. Had our implementation been slightly
different---had we, say, cached the acknowledgement path in a way that
happened to refuse---the test would have passed, and we would have shipped a
green fail-closed test over a durability path that had never once been made to
fail.

That is the shape that worries me. A test that passes for the wrong reason is
worse than a missing test, because the missing test is visible in a coverage
report and the wrong one is not. And this particular wrong reason is
attractive: ``delete the file'' is the first thing anybody tries.

\subsection{How we actually found it, which was luck}

I want to be accurate about the process, because the tidy version would be
dishonest.

We did not find this by reasoning about POSIX. We found it because the
assertion failed---and it failed in the direction that says ``you acknowledged
something you should have refused,'' which is alarming enough that we went
looking at the mechanism instead of the implementation. If the test had passed,
we would have ticked it off and moved on. The mechanism was only inspected
because the outcome was already wrong.

So the honest summary is: a bug in our fault injection was caught by an
unrelated bug in the same test, and if either had been absent we would have
shipped. We wrote a standalone probe afterwards to confirm the POSIX behaviour
directly rather than infer it from the failure, and that probe is the listing
above.

\subsection{The fix, and what it does not fix}

The repair is not in the test. A committer that cannot re-open its store now
refuses everything, because a descriptor that has outlived its path is exactly
the condition under which acknowledging is unsafe. That is production
behaviour, not a test affordance: the same reasoning applies if an operator
moves a log directory under a running process.

What we have not done is establish that our durability path is correct under
real crash conditions. Refusing when the store vanishes is one narrow case.
The literature on crash-consistency testing \cite{pillai2014,zheng2014} exists
because doing this properly means enumerating reorderings a real storage stack
permits, and we have not done that. We tested that a torn trailing record is
truncated and that a corrupted complete record is refused, by writing the
damage ourselves. That is a simulation of a crash, not a crash.

\section{How We Checked}
\label{sec:method}

The proofs are the part of this system where a subtle error is both easy to
make and hard to notice, so they get the disproportionate share of the
testing, and the testing is exhaustive rather than sampled.

For every tree size from 1 to 33, and for every leaf in the tree, we generate
an inclusion proof and verify it. For every tree size from 1 to 33, and every
prefix length from 1 to that size, we generate a consistency proof and verify
it. Separately, the incremental root---maintained as a right fringe as leaves
arrive---is compared against a recursive reference implementation written
straight from the RFC, at every size.

The bound of 33 is chosen to get past the powers of two. RFC~6962 splits every
interior node at the largest power of two strictly below the node's size, and
an implementation that gets this subtly wrong is usually still correct at 2, 4,
8, 16 and 32. A test suite that checks those sizes finds nothing. Sizes 17 and
33 are where the fringe logic has to do something interesting, which is why the
range ends just past one.

The discipline we set ourselves, and kept, was: if a proof test fails, fix the
verifier against the RFC, never relax the test. This sounds obvious written
down. It is less obvious at 2am when the exhaustive test fails at size 23 and
the tempting move is to check a narrower property. We record it here because
the sibling project's history contains exactly that failure: a verifier that
compared record counts instead of roots, and consequently accepted a rewritten
history, which is the attack in Section~\ref{sec:problem} passing a test suite.

We also assert the negative controls, which matter more than they look. It is
easy to write a verifier that rejects a rewritten log; the empty verifier that
rejects everything does that. So there is a test that an \emph{honest}
extension of the same prior root is accepted, and a test asserting that a
rewritten log is accepted when no prior root is supplied---because that
acceptance is the honest limitation, and if it ever changed silently, the tool
would be claiming a defence it does not have.

\section{The Chain That Does Not Cover The Record}
\label{sec:gap}

One finding about the specification we were implementing, which is smaller
than Section~\ref{sec:test} but easier to act on.

The evidence schema defines the chain head as

\[
h_{i+1} = H(h_i \; \| \; \texttt{canonical\_snapshot\_hash}_i \; \| \; \texttt{verdict}_i)
\]

that is, over two fields of the record plus the previous head. Not over the
record.

So the chain protects those two fields. Everything else in the record is
outside it: the step index, the target resource, the policy bundle hash, the
agent identifier, the initiating user, the timestamp. An operator---or anyone
who can write to the store---can change which resource an action touched, or
which user initiated it, or which policy bundle was in force, and every chain
link still verifies, because none of those values is an input to any link.

That is very likely an oversight rather than a decision; a snapshot hash and a
verdict are the two fields that feel like ``the decision,'' and it is easy to
write the definition around them. But the consequence is that a log which
passes its own chain check has not established the integrity of most of its
content.

We chain over the whole record digest instead:
$h_{i+1} = H(h_i \| d(r_i))$. It is strictly stronger, and it is free---$d(r_i)$
is already computed to form the Merkle leaf, so this costs one hash we were
performing anyway. We report the schema definition upstream as a finding rather
than implement it. An implementation that honoured it would be conformant and
weaker, and we would rather be non-conformant in the direction of covering more
of the record.

\section{What We Still Don't Know}
\label{sec:unknown}

The single largest gap is the one this paper is about, and we have not closed
it: we have no publication point. Everything here reduces to the custody
condition, and the custody condition is an arrangement between organisations
that a Rust crate cannot create. What the crate can do---and does---is refuse
to let a root be printed without the sentence saying so, so that nobody reads a
green tick as more than it is. That is a documentation defence against a
structural problem, and I would not want to oversell it.

We have no performance numbers. None. The log batches claims into a single
sync and we have not measured what that costs, so we do not know whether the
durability requirement is compatible with the action rates people want from
agents, and we are not going to guess. Anyone who needs that number should
treat this paper as silent on it rather than as reassuring.

We have not measured the linkability cost of the tree topology. A single global
tree reveals nothing about who acted, because membership distinguishes nobody.
A tree per agent parallelises well and its key \emph{is} an agent identifier,
so every action in it is linked to every other by construction. Those two facts
point in opposite directions and we have implemented all three arrangements
without measuring either axis.

The signer is a software key, and we should be blunt that this means the whole
system currently has no defence against its operator at all---the custody
argument above is not mitigated by the signer, it is the reason the signer's
weakness does not much matter yet. A key in a trusted execution environment
would change who can produce a signature, not who can publish a root, so it
narrows the attack without closing it.

Finally, on the finding in Section~\ref{sec:test}: we do not know how common
it is. We found one instance, in our own code, by accident. The reason we
think it generalises is that the reasoning that produced our broken test---``to
make a resource unavailable, delete it''---is not idiosyncratic, and the
failure is silent. But we have not surveyed other durability test suites for
the same pattern, and until somebody does, ``we think this is common'' is a
hypothesis and should be read as one.

\section*{Availability}

The implementation, the runnable demonstration of the rewrite attack, and the
source of this paper are at
\url{https://github.com/Task-force-for-AI-agents-in-Healthcare/ephemeris},
under Apache-2.0.

\bibliographystyle{plain}
\bibliography{references}

\end{document}
